%%% LaTeX-Vorlage Version 1.8 %%%

% Grundlegende Dokumenteneigenschaften gemäß DHBW-Vorgaben
\documentclass[a4paper,fontsize=11pt,oneside,parskip=half,headings=normal]{scrreprt} 
% \usepackage{showframe} % nur für Kontrolle der Ränder 

%%% Präambel einbinden (mit Festlegungen gemäß DHBW-Vorgaben) %%%
%%% Präambel %%%
% hier sollten keine Änderungen erforderlich sein
%
\usepackage[utf8]{inputenc}   % Zeichencodierung UTF-8 für Eingabe-Dateien
\usepackage[T1]{fontenc}      % Darstellung von Umlauten im PDF

\usepackage{listings}         % für Einbindung von Code-Listings
\lstset{numbers=left,numberstyle=\tiny,numbersep=5pt,texcl=true}
\lstset{literate=             % erlaubt Sonderzeichen in Code-Listings 
{Ö}{{\"O}}1
{Ä}{{\"A}}1
{Ü}{{\"U}}1
{ß}{{\ss}}2
{ü}{{\"u}}1
{ä}{{\"a}}1
{ö}{{\"o}}1
{€}{{\euro}}1
}

\usepackage[
  inner=35mm,outer=15mm,top=25mm,
  bottom=20mm,foot=12mm,includefoot
]{geometry}                 % Einstellungen für Ränder

\usepackage[english]{babel} % Spracheinstellungen Englisch
\usepackage[babel,english=british]{csquotes} % englische Anf.zeichen
\usepackage{enumerate}      % anpassbare Nummerier./Aufz.
\usepackage{graphicx}       % Einbinden von Grafiken
\usepackage[onehalfspacing]{setspace} % anderthalbzeilig

\usepackage{blindtext}      % Textgenerierung für Testzwecke
\usepackage{color}          % Verwendung von Farbe 

\usepackage{acronym}        % für ein Abkürzungsverzeichnis

\usepackage[                % Biblatex
  backend=biber,
  bibstyle=_dhbw_authoryear,maxbibnames=99,
  citestyle=authoryear,     
  uniquename=true, useprefix=true,
  bibencoding=utf8]{biblatex}
%kein Punkt am Ende bei \footcite
%http://www.golatex.de/footcite-ohne-punkt-am-schluss-t4865.html
\renewcommand{\bibfootnotewrapper}[1]{\bibsentence#1}


%Reihenfolge der Autorennamen
%   
% http://golatex.de/viewtopic,p,80448.html#80448
% Argumente: siehe http://texwelt.de/blog/modifizieren-eines-biblatex-stils/
\DeclareNameFormat{sortname}{% Bibliographie
  \ifnum\value{uniquename}=0 % Normalfall
    \ifuseprefix%
      {%
         \usebibmacro{name:family-given}
           {\namepartfamily}
           {\namepartgiveni}
           {\namepartprefix}
           {\namepartsuffixi}%
       }
      {%
         \usebibmacro{name:family-given}
           {\namepartfamily}
           {\namepartgiveni}
           {\namepartprefixi}
           {\namepartsuffixi}%
       }%
  \fi
  \ifnum\value{uniquename}=1% falls nicht eindeutig, abgek. Vorname 
      {%
         \usebibmacro{name:family-given}
           {\namepartfamily}
           {\namepartgiveni}
           {\namepartprefix}
           {\namepartsuffix}%
       }%
  \fi
  \ifnum\value{uniquename}=2% falls nicht eindeutig, ganzer Vorname 
      {%
         \usebibmacro{name:family-given}
           {\namepartfamily}
           {\namepartgiven}
           {\namepartprefix}
           {\namepartsuffix}%
       }%
  \fi   
  \usebibmacro{name:andothers}}

\DeclareNameFormat{labelname}{% für Zitate
  \ifnum\value{uniquename}=0 % Normalfall
    \ifuseprefix%
      {%
         \usebibmacro{name:family-given}
           {\namepartfamily}
           {\empty}
           {\namepartprefix}
           {\namepartsuffixi}%
       }
      {%
         \usebibmacro{name:family-given}
           {\namepartfamily}
           {\empty}
           {\namepartprefixi}
           {\namepartsuffixi}%
       }%
  \fi
  \ifnum\value{uniquename}=1% falls nicht eindeutig, abgek. Vorname 
      {%
         \usebibmacro{name:family-given}
           {\namepartfamily}
           {\namepartgiveni}
           {\namepartprefix}
           {\namepartsuffix}%
       }%
  \fi
  \ifnum\value{uniquename}=2% falls nicht eindeutig, ganzer Vorname 
      {%
         \usebibmacro{name:family-given}
           {\namepartfamily}
           {\namepartgiven}
           {\namepartprefix}
           {\namepartsuffix}%
       }%
  \fi   
  \usebibmacro{name:andothers}}
      
  
\DeclareFieldFormat{extrayear}{% = the 'a' in 'Jones 1995a'
  \iffieldnums{labelyear}
    {\mknumalph{#1}}
    {\mknumalph{#1}}}        

\renewcommand*{\multinamedelim}{\addslash}
\renewcommand*{\finalnamedelim}{\addslash}
\renewcommand*{\multilistdelim}{\addslash}
\renewcommand*{\finallistdelim}{\addslash}

\renewcommand{\nameyeardelim}{~}

% Literaturverzeichnis: Doppelpunkt zwischen Name (Jahr): Rest 
% http://de.comp.text.tex.narkive.com/Tn1HUIXB/biblatex-authoryear-und-doppelpunkt
\renewcommand{\labelnamepunct}{\addcolon\addspace}

% damit die Darstellung für Vollzitate von Primärquellen in 
% Fußnoten später auf "nicht fett" geändert werden kann 
% (nur für Zitate von Sekundärliteratur relevant)
\newcommand{\textfett}[1]{\textbf{#1}}

% für Zitate von Sekundärliteratur:
\newcommand{\footcitePrimaerSekundaer}[4]{%
  \renewcommand{\textfett}[1]{##1}%
  \footnote{\fullcite[#2]{#1}, zitiert nach \cite[#4]{#3}}%  
  \renewcommand{\textfett}[1]{\textbf{##1}}%
}

% Im Literaturverzeichnis: Autor (Jahr) fett
\renewbibmacro*{author}{%
  \ifboolexpr{%
    test \ifuseauthor%
    and
    not test {\ifnameundef{author}}
  }
    {\usebibmacro{bbx:dashcheck}
       {\bibnamedash}
       {\usebibmacro{bbx:savehash}%
        \textfett{\printnames{author}}%
        \iffieldundef{authortype}
          {\setunit{\addspace}}
          {\setunit{\addcomma\space}}}%
     \iffieldundef{authortype}
       {}
       {\usebibmacro{authorstrg}%
        \setunit{\addspace}}}%
    {\global\undef\bbx@lasthash
     \usebibmacro{labeltitle}%
     \setunit*{\addspace}}%
  \textfett{\usebibmacro{date+extrayear}}}

% Sonderfall: Quelle ohne Autor, aber mit Herausgeber
% Name des Herausgebers wird fett gedruckt
\renewbibmacro*{bbx:editor}[1]{%
  \ifboolexpr{%
    test \ifuseeditor%
    and
    not test {\ifnameundef{editor}}
  }
    {\usebibmacro{bbx:dashcheck}
       {\bibnamedash}
       {\textfett{\printnames{editor}}%
        \setunit{\addcomma\space}%
        \usebibmacro{bbx:savehash}}%
     \usebibmacro{#1}%
     \clearname{editor}%
     \setunit{\addspace}}%
    {\global\undef\bbx@lasthash
     \usebibmacro{labeltitle}%
     \setunit*{\addspace}}%
  \textfett{\usebibmacro{date+extrayear}}}

% Anpassungen für deutsche Sprache
\DefineBibliographyStrings{english}{%
	nodate = {{o.J.}},
	urlseen = {{Abruf:}},
	ibidem = {{ebenda}}
}

% keine Anführungszeichen beim Titel im Literaturverzeichnis
\DeclareFieldFormat[article,book,inbook,inproceedings,manual,misc,phdthesis,thesis,online,report]{title}{#1\isdot}

\newcommand{\literaturverzeichnis}{%
% nur Literaturverzeichnis
% (als eigenes Kapitel)
\phantomsection
\addcontentsline{toc}{chapter}{List of literature}
\spezialkopfzeile{List of literature}
\defbibheading{lit}{\section*{List of literature}}
\label{chapter:quellen}
\printbibliography[heading=lit,notkeyword=ausblenden]
}

% Anpassungen für englische Sprache
\DefineBibliographyStrings{english}{
    nodate = {{w.y.}},
    urlseen = {{retrieval:}}
} % mit DHBW-spezifischen Einstellungen

\usepackage{hyperref}       % URL-Formatierung, klickbare Verweise

\usepackage{tocloft}        % für Verzeichnis der Anhänge

\usepackage{graphicx}       % Einbinden von Grafiken
\usepackage{tikz}           %  -"-

\newcounter{anhcnt}
\setcounter{anhcnt}{0}
\newlistof{anhang}{app}{}

\newcommand{\anhang}[1]{%
  \refstepcounter{anhcnt}
  \setcounter{anhteilcnt}{0}
  \section*{Appendix \theanhcnt: #1}
  \addcontentsline{app}{section}{\protect\numberline{Anhang \theanhcnt}#1}\par
}

\newcounter{anhteilcnt}
\setcounter{anhteilcnt}{0}

\newcommand{\comment}[1]{} % Ermoeglicht das einfuegen von commentaren, die nicht gerendert werden

\newcommand{\anhangteil}[1]{%
	\refstepcounter{anhteilcnt}
	\subsection*{Anhang~\arabic{anhcnt}/\arabic{anhteilcnt}: #1}
	\addcontentsline{app}{subsection}{\protect\numberline{Anhang \theanhcnt/\arabic{anhteilcnt}}#1}\par
}

\renewcommand{\theanhteilcnt}{Anhang \theanhcnt/\arabic{anhteilcnt}}

% vgl. S. 4 Paket-Beschreibung tocloft 	
% Einrückungen für Anhangverzeichnis
\makeatletter
\newcommand{\abstaendeanhangverzeichnis}{
\renewcommand*{\l@section}{\@dottedtocline{1}{0em}{5.5em}}
\renewcommand*{\l@subsection}{\@dottedtocline{2}{2.3em}{6.5em}}
}
\makeatother

% Abbildungs- und Tabellenverzeichnis
% Bezeichnungen
\renewcaptionname{english}{\figurename}{Abb.}
\renewcaptionname{english}{\tablename}{Tab.}
% Einrückungen
\makeatletter
\renewcommand*{\l@figure}{\@dottedtocline{1}{0em}{2.3em}}
\renewcommand*{\l@table}{\@dottedtocline{1}{0em}{2.3em}}
\makeatother


\usepackage{chngcntr}                % fortlaufende Zähler für Fußnoten, Abbildungen und Tabellen
\counterwithout{figure}{chapter}
\counterwithout{table}{chapter}
\counterwithout{footnote}{chapter}

\usepackage[automark]{scrlayer-scrpage} 
%% Definitionen für Kopf- und Fußzeile auf normalen Seiten
\defpagestyle{kopfzeile}
{% Kopfdefinition
  (\textwidth,0pt)    % Länge der oberen Linie,Dicke der oberen Linie       
  {} % Definition für linke Seiten im doppelseitigen Layout
  {} % Definition für rechte Seiten im doppelseitigen Layout      
  {  % Definition für Seiten im einseitigen Layout
	\makebox[0pt][l]{\rightmark}% 
	\makebox[\linewidth]{}% 
  }        
  (\textwidth, 0.4pt) % Untere Linienlänge, Untere Liniendicke
}
{% Fußdefinition
  (\textwidth,0pt)    % Obere Linienlänge, Obere Liniendicke
  {} % Definition für linke Seiten im doppelseitigen Layout
  {} % Definition für rechte Seiten im doppelseitigen Layout
  {  % Definition für Seiten im einseitigen Layout
    \makebox[\linewidth]{}%
    \makebox[0pt][r]{\pagemark}%
  }
  (\textwidth, 0pt)   % Länge der unteren Linie,Dicke der unteren Linie
}

%% Definitionen für Kopf- und Fußzeile auf ersten Seiten eines Kapitels
\defpagestyle{kapitelkopfzeile}
{% Kopfdefinition
  (\textwidth,0pt)    % Länge der oberen Linie,Dicke der oberen Linie       
  {} % Definition für linke Seiten im doppelseitigen Layout
  {} % Definition für rechte Seiten im doppelseitigen Layout      
  {}  % Definition für Seiten im einseitigen Layout
  (\textwidth, 0pt) % Untere Linienlänge, Untere Liniendicke
}
{% Fußdefinition
  (\textwidth,0pt)    % Obere Linienlänge, Obere Liniendicke
  {} % Definition für linke Seiten im doppelseitigen Layout
  {} % Definition für rechte Seiten im doppelseitigen Layout
  {  % Definition für Seiten im einseitigen Layout
    \makebox[\linewidth]{}%
    \makebox[0pt][r]{\pagemark}%
  }
  (\textwidth, 0pt)   % Länge der unteren Linie,Dicke der unteren Linie
}

%% Definitionen für Kopf- und Fußzeile im Anhang und bei Quellenverzeichnisse
\newcommand{\spezialkopfzeileBezeichnung}{}
\defpagestyle{spezialkopfzeile}
{% Kopfdefinition
  (\textwidth,0pt)    % Länge der oberen Linie,Dicke der oberen Linie       
  {} % Definition für linke Seiten im doppelseitigen Layout
  {} % Definition für rechte Seiten im doppelseitigen Layout      
  {  % Definition für Seiten im einseitigen Layout
	\makebox[0pt][l]{\spezialkopfzeileBezeichnung}% 
	\makebox[\linewidth]{}% 
  }        
  (\textwidth, 0.4pt) % Untere Linienlänge, Untere Liniendicke
}
{% Fußdefinition
  (\textwidth,0pt)    % Obere Linienlänge, Obere Liniendicke
  {} % Definition für linke Seiten im doppelseitigen Layout
  {} % Definition für rechte Seiten im doppelseitigen Layout
  {  % Definition für Seiten im einseitigen Layout
    \makebox[\linewidth]{}%
    \makebox[0pt][r]{\pagemark}%
  }
  (\textwidth, 0pt)   % Länge der unteren Linie,Dicke der unteren Linie
}
            
\newcommand\spezialkopfzeile[1]{%
  \renewcommand\spezialkopfzeileBezeichnung{#1}
  \pagestyle{spezialkopfzeile}
}
                
% Standard-Pagestyle auswählen
\pagestyle{kopfzeile}

% keine Kopfzeile anzeigen auf Seiten, auf denen ein 
% Kapitel beginnt oder das Inhalts-/Abbildungs-/Tabellenverzeichnis steht 
\renewcommand{\chapterpagestyle}{kapitelkopfzeile}
\tocloftpagestyle{kapitelkopfzeile}

		 % für schöne Kopfzeilen 

\usepackage{textcomp}            % erlaubt EUR-Zeichen in Eingabedatei
\usepackage{eurosym}             % offizielles EUR-Symbol in Ausgabe
\renewcommand{\texteuro}{\euro}  % ACHTUNG: nach hyperref aufrufen!

\usepackage{scrhack}             % stellt Kompatibilität zw. KOMA-Script
                                 % (scrreprt) und anderen Paketen her
                                 
% Anpassung der Abstände bei Kapitelüberschriften
% (betrifft auch Inhalts-, Abbildungs- und Tabellenverzeichnis)
\renewcommand*\chapterheadstartvskip{\vspace*{-\topskip}}
\newcommand{\myBeforeTitleSkip}{1mm}
\newcommand{\myAfterTitleSkip}{10mm}
\setlength\cftbeforetoctitleskip{\myBeforeTitleSkip}
\setlength\cftbeforeloftitleskip{\myBeforeTitleSkip}
\setlength\cftbeforelottitleskip{\myBeforeTitleSkip}

\setlength\cftaftertoctitleskip{\myAfterTitleSkip}
\setlength\cftafterloftitleskip{\myAfterTitleSkip}
\setlength\cftafterlottitleskip{\myAfterTitleSkip}                                                            
%%% Ende der Präambel %%%

%%% Name der eigenen Literatur-Datenbank (ggf. anpassen) %%%
\bibliography{includes/boilerplate/literatur-datenbank.bib}

\begin{document}
%%% Deckblatt einbinden %%% 
% Anpassungen nötig (Name, Titel etc.)
% HIER EDITIEREN: 
% Typ der Arbeit (für Deckblatt und ehrenwörtliche Erklärung)
% - bitte Zutreffendes auswählen
\newcommand{\typMeinerArbeit}{2. Project Paper} 


% Thema der Arbeit (für ehrenwörtliche Erklärung, ohne Umbrüche)
% HIER EDITIEREN: 
\newcommand{\themaMeinerArbeit}{"Development and demonstration of a proof-of-concept for the integration of programming frameworks for high performance computing into a container-based workflow orchestrator."}

% Vorname, Name der Autorin/des Autors (für Titelseite und Metadaten)
% HIER EDITIEREN:
\newcommand{\meinName}{Jon Eckerth}

\thispagestyle{empty}

% Include DHBW and Corporate Logo

\begin{center}	
\begin{tikzpicture}:
    \node[midway,xshift=-5cm](c){\includegraphics[width=0.35\textwidth]{graphics/Hewlett_Packard_Labs_logo.png}};
    \node[midway,xshift=5cm](c){\includegraphics[width=0.35\textwidth]{graphics/DHBW_Logo.png}};
\end{tikzpicture}
\end{center}


\begin{spacing}{1}
\begin{center}	
~\vspace{0mm}

% HIER EDITIEREN: Titel der Arbeit
{\sffamily
\large  
% \Large  % bei sehr langen Titeln ggf. etwas kleinere Schriftart wählen
\textbf{
    Development and demonstration of a proof-of-concept for the integration of programming frameworks for high performance computing into a container-based workflow orchestrator.
    }
\bigskip
\textbf{}
}


\vspace{15mm}

% Typ wird automatisch eingefügt (oben festlegen)
{\Large \typMeinerArbeit}

\vspace{1cm}

% HIER ggf. EDITIEREN
submitted on \today 

\vspace{15mm}

Faculty for Business and Health
\medskip

%Studiengang:
Degree program in Business Informatics

\medskip

% HIER EDITIEREN: Kurs eintragen
International Management for Business and Information Technology 

\vspace{10mm}

by

\vspace{10mm}
    
% Vorname und Name wird automatisch eingefügt (oben festlegen) 
{\large\textsc{\meinName}}

\vspace{10mm}
\end{center}

\vfill

% HIER EDITIEREN: Name des Unternehmens, Name der Betreuerin/des Betreuers
\begin{tabular}{ll}
Practice supervisor : & DHBW Stuttgart: \\
\hspace{0.4\linewidth} & \\
$\langle$ Hewlett Packard Labs $\rangle$ & $\langle$ Dominic Viola $\rangle$ \\
$\langle$ Harumi Kuno, PhD $\rangle$ 
&  \\
$\langle$ Principal Research Scientist $\rangle$ \\
\\
Signature of the practice supervisor: \\
\end{tabular}


\vspace{1cm}
%(etwas Platz für die Unterschrift der Betreuerin/des Betreuers aus der Ausbildungsstätte)
\end{spacing}

% falls ein Vertraulichkeitsvermerk erforderlich ist,
% die Kommentarzeichen in den nachfolgenden Zeilen entfernen:
 
% \begin{center}
% \small
% \textbf{Vertraulichkeitsvermerk}:
% Der Inhalt dieser Arbeit darf weder als Ganzes noch in Auszügen \\
% Personen außerhalb des Prüfungs- und Evaluationsverfahrens zugänglich gemacht werden, sofern keine anders lautende Genehmigung des Dualen Partners vorliegt. 
% \end{center}

% Meta-Daten für PDF-Datei basierend auf obigen Angaben
\hypersetup{pdftitle={\themaMeinerArbeit}}
\hypersetup{pdfauthor={\meinName}}
\hypersetup{pdfsubject={\typMeinerArbeit\ DHBW Stuttgart \the\year}}

%%% Declaration of originality einbinden %%%
\clearpage
\thispagestyle{empty}

{\LARGE\textsf{\textbf{ Author’s declaration}}}

% \typMeinerArbeit und \themaMeinerArbeit werden in deckblatt.tex definiert

Hereby I solemnly declare:
\begin{enumerate}
    \item that this \typMeinerArbeit\, titled \emph{\themaMeinerArbeit} is entirely the product of my own scholarly work, unless otherwise
    indicated in the text or references, or acknowledged below;

    \item I have indicated the thoughts adopted directly or indirectly from other sources at
    the appropriate places within the document;

    \item this \typMeinerArbeit\ has not been submitted either in whole or
    part, for a degree at this or any other university or institution;

    \item  I have not published this \typMeinerArbeit\ in the past;

    \item the printed version is equivalent to the submitted electronic one.
\end{enumerate}
I am aware that a dishonest declaration will entail legal consequences.


\vspace{2cm}

Dublin, \today

\begin{tikzpicture}
\draw[thick](0,0)--(5,0)node[midway,yshift=0.3cm](c){\includegraphics[width=4cm]{graphics/place_holder_signature.png}};
%\draw[thick](0,0)--(5,0)node[midway,yshift=0.3cm](c){\includegraphics[width=5cm]{graphics/signature.png}};

\node[yshift=-2mm] at (c.south) {\bfseries \meinName};
\end{tikzpicture}


%%% Umstellung der Seiten-Nummerierung auf i, ii, iii ... %%%
\pagenumbering{Roman} 

%%% Abstract einbinden (optionale Kurzfassung Ihrer Arbeit) %%%
\begin{abstract}
\thispagestyle{kapitelkopfzeile}
\textbf{Abstract:}

This Second project paper presents the development and demonstration of a proof-of-concept for the
integration of \ac{HPC} programming frameworks into a container-base workflow orchestrator.

This project aim to reconcile the different demands of the \ac{HPC} and \ac{CC} communities by demonstrating the integration of
the \ac{HPC} programming framework Arkouda into the container-based workflow orchestrator Pachyderm, showing the technical feasibility of this approach.

A prototype implementing this integrated system is constructed and evaluated through prototyping methodologies, with a focus on resilience, scalability, portability, and user-friendliness.
The prototype is iteratively refined to address \ac{LCP} and \ac{TCP}, with particular attention to the usability of the system for non \ac{CC} experts.

This project paper contributes to the body of knowledge by way of practical example, lessons learned with each iteration, and sheds light on pathways for future research towards
a landscape where the seamless and efficient integration of \ac{HPC} workloads in \ac{CC} environments becomes possible

\end{abstract}
\cleardoublepage

%%% Inhalts-, Abbildungs-, Tabellenverzeichnisse %%%
% sollen einzeilig gesetzt werden, um Platz zu sparen 
\begin{spacing}{1}
\tableofcontents
\clearpage
\chapter*{List of abbreviations}
\addcontentsline{toc}{chapter}{List of abbreviations}

\begin{acronym}[DHBW] 
% Argument definiert die Breite der ersten Spalte anhand des längsten vorkommenden Eintrags
\acro{AI}{Artificial Intelligence}
\acro{k8s}{Kubernetes}
\acro{ML}{Machine Learning}
\acro{HPC}{High Performance Computing}
\acro{CC}{Cloud Computing}
\acro{CI/CD}{Continuous Integration/Continuous Delivery}
\acro{IaC}{Infrastructure as Code}
\acro{PoC}{Proof of Concept}
\acro{HPE}{Hewlett Packard Enterprise}
\acro{CNCF}{Cloud Native Computing Foundation}
\acro{CLASP}{Cloud Application Services Platform}
\acro{CWL}{Common Workflow Language}
\acro{SME}{Subject Matter Expert}
\acro{TCPP}{Tightly Coupled Parallel Problem}
\end{acronym}



\clearpage
\thispagestyle{kapitelkopfzeile}
\listoffigures
\phantomsection
\addcontentsline{toc}{chapter}{List of figures} % Abb.verz. ins Inh.verz. aufnehmen

\clearpage
\listoftables
\phantomsection
\addcontentsline{toc}{chapter}{List of tables}    % Tab.verz. ins Inh.verz. aufnehmen
\end{spacing}    

%%% Umstellung der Seiten-Nummerierung auf 1, 2, 3 ... %%%
\cleardoublepage
\pagenumbering{arabic}


% ----------------------- Inhalt -------------------------------------


\chapter{Introduction}
\label{Introduction}

In this section, the underlying motivation of this project is explained.
Furthermore, the problems which will be addressed by this project are described,
which serve as the basis for the research questions which will guide this project and ultimately
result in solutions and further questions which are listed in the contributions section and discussed in the conclusion.

\section{Motivation}

The proliferation of "Big Data" has led to the need to compute, analyse, and visualize ever-increasing amounts of datasets,
which themselves are getting more and more complex, has led to an ever increasing demand for more efficient and quicker ways to process data.

Both the \ac{HPC} and the \ac{CC} community have been working on solutions to distribute and parallelize computations for decades, 
both with their own approaches and solutions to their respective problems.

While the \ac{HPC} community has been putting a lot of effort into developing new and extremely efficient ways to parallelize computations,
the \ac{CC} community has been focusing on improving the flexibility, scalability and resilience of their solutions as well as improving the ease of use for their developers and users. 

Both used to be very distinct and separate communities due to their very different usecases, while the \ac{HPC} community was mostly concerned with scientific computing and simulations of physical phenomena,
the \ac{CC} community is mostly concerned with providing a reliable and easily up and down scalable infrastructure for the industry and businesses.

Now with the advent of \ac{ML} and \ac{AI} the two communities are starting to converge, 
as the \ac{ML} and \ac{AI} community is adopting the tools and techniques of both communities to solve their problems as they see fit.

But this convergence of the two is not without its problems, being developed in two coexisting and separate communities, the tools and techniques of both communities are not always compatible with each other,
the goal of this project is to find a way to bridge this gap and to find a way to combine the best of both worlds.

\newpage

\section{Problem Statement}

The following key problems have emerged from the convergence of High Performance Computing (HPC) and Cloud Computing (CC) communities, especially in the context of Machine Learning (ML) and Artificial Intelligence (AI) research:

\begin{itemize}

    \item \textbf{Workload Resilience and Fault Tolerance in \ac{HPC}:} 
        \ac{HPC} systems often lack mechanisms to recover from task failures within
        larger jobs, running for an extended time.
        When a component task fails, it can invalidate the entire
        computation, requiring a restart from scratch.
        This need for resilient failover and verification strategies as well as the need to avoid
        computational wastage is a key challenge for \ac{HPC} systems, especially with every increasing system sizes and complexity.
        \footcite{egwutuohaSurveyFaultTolerance2013}

    \item \textbf{Environment/Package Management in \ac{HPC}:} 
        \ac{HPC} systems are notorious for their complex package management systems.
        As having a shared infrastructure between many users each with their own specific needs and requirements of different versions of packages, libaries and software,
        all the while sharing a common environment.
        Many solutions to this problem have been developed, each with their own advantages and disadvantages.
        \footcite{duboisWhyJohnnyCan2003} \footcite{bzeznikNixHPCPackage2017} \footcite{gamblinSpackPackageManager2015} \footcite{hosteEasyBuildBuildingSoftware2012}

    \item \textbf{Portability Issues with \ac{HPC}:}
        Tieing in with the previous point, \ac{HPC} systems are often designed to be optimized for specific hardware as well as having a very specific software stack.
        This makes the portability of applications between different \ac{HPC} systems very difficult and often infeasible. \footcite[p. 50]{canonCasePortabilityReproducibility2019}
        This lack of portability contrasts sharply with the more platform-agnostic nature of \ac{CC} environments, where the containerization of applications has become the norm for ensuring portability.

    \item \textbf{Scalability and Flexibility in \ac{HPC}:}
        Due to its direct access to the hardware and very specific hardware needs, \ac{HPC} systems are often hard to dynamically scale and inflexible. 
        while \ac{CC} systems are designed to be easily scalable and flexible and are often designed to be hardware agnostic and abstract away the underlying hardware.
        This becomes especially relevant in the context of heterogeneous hardware, where the hardware is not uniform and consists of different types of hardware,
        which is becoming more and more common in the context of ML and AI research.
      
    \item \textbf{Lack of Interconnected Problem Solving in \ac{CC}:} 
        The workloads traditionally deployed on \ac{CC} systems are often independent of each other, like load balancing, web hosting, etc. 
        This is in stark contrast to the interconnected nature of \ac{HPC} workloads, where each part of the input data is dependent on the other parts of the input data,
        such that all nodes of the system need to be able to communicate with each other.

    \item \textbf{Provenance and Reproducibility:}
        Another need that is becoming more and more important in the context of ML and AI research is the need for provenance and reproducibility of results.
        Being able to tell which data was used to train the model, is of ever increasing importance as the influence the resulting models have on our lives increases as well as the data used to train the model.
        This is especially important since it is crucial to ensure that the data is not biased, outdated, or otherwise flawed, which could lead to incorrect predictions, decisions, or recommendations. 
        in addition various data sources, from images to text, may have copyright restrictions that, when overlooked, can lead to significant legal complications.

    \item \textbf{Versioning Limitations:} 
        The dynamic nature of ML and AI research necessitates robust versioning solutions for data, configurations and code.
        While \ac{CC} has developed many solutions to this problem over the years, making them their own subsection of the ecosystem,
        namely \ac{CI/CD} tools for the testing and deployment of applications aswell as \ac{IaC} tools for the deployment of infrastructure.
        While many solutions have been developed for the one-off deployment of \ac{HPC} systems, the dynamic nature of \ac{CC} systems necessitates a more robust solution to this problem, 
        from which the \ac{HPC} community could benefit aswell. 
        
\end{itemize}

\section{Research Questions}

To address the aforementioned problems, to bridge the gap between the two paradigms and to combine the best of both worlds, an integration of the two paradigms is needed.
This was accomplished by integrating a \ac{HPC} framework called 'Arkouda \footcite{merrillArkoudaInteractiveData2019}' into a container based \ac{CC} workflow management tool called 'Pachyderm'\footcite{pachydermPachyderm} and integrating both with the supporting infrastructure the \ac{CC} system enables us to use.
This process of integration and prototyping as well as the explanation of the underlying concepts and technologies will be the focus of this project.

\begin{itemize}
    \item \textbf{RQ1:} \textit{
        How can a high-performance computing framework be effectively integrated into a container-based workflow management tool?
    }
    \item \textbf{RQ2:} \textit{
        What are the opportunities for improving the integration of high-performance computing frameworks with container-based workflow management tools?
    }
\end{itemize}

\newpage

\section{Contributions}

In order to address the problems stated above, find answers to the research questions and to bridge the gap between the two paradigms, the following contributions were made:
\begin{itemize}
    \item \textbf{C1:} \textit{ An analysis of the problem space and existing solution, within the constraints of time, resources and businesses needs.}
    \item \textbf{C2:} \textit{ A prototype implementation combining the 'Arkouda' framework with the \ac{k8s} based workflow orchestrator 'Pachyderm'}
    \item \textbf{C3:} \textit{ Further integrations of tools from both sides of the spectrum, addressing many of aforementioned pain-points}
\end{itemize}

\chapter{Methodology}
\label{methodology}

\section{Prototyping}


The methodology of Prototyping in the context of software development has been a well established practice for many years and
has been discussed in literature for at least the last 40 years \footcite{gomaaImpactRapidPrototyping1983}.
It is especially useful in complex situations where the requirements might not be fully known or understood at the beginning of the project
and need to be discovered in an iterative and exploratory process, 
as it enables the project to be partitioned into more manageable chunks which can be tackled individually\footcite{naumannPrototypingNewParadigm1982} .

At their core prototyping is a state transition model, where a finial state is approached in a series of non-ideal intermediate steps,
each of which revealing new information about the problem domain and the solution space and applying this new information to the next iteration\footcite{kraushaarPrototypingMethodApplications1985}.

It is described as the methodology of with the highest constructive plasticity, as it supports many different flavors and variations of the core concept,
but as long as it follows the core principles of this iterative approximation of an ideal version, and specifically evaluates the results of the iterations it should be considered a methodologically complete scientific method\footcite{wildeMethodenspektrumWirtschaftsinformatikUeberblick} .

For this project we will be using a concatenated version of the so called \ac{RAD} model, which has been quite established in software development, as it is very well suited for projects with a high degree of complexity and uncertainty, while reducing the admistrative overhead of the project\footcite{martinRapidApplicationDevelopment1991}.
Reducing the administrative overhead is especially important for this project, as the time frame is limited to 3 months, and the project is being done by a single person, 
therefore more complex tools like the spiral model \footcite{boehmSpiralModelSoftware1988}, which is a popular model emphasizing a risk analysis with each step of the iteration, 
would be too time-consuming to be feasible for this project.

First there will be a central initial goal which will be the goal of what we are trying to reach,
as the requirements will change with the knowledge acquired and the scope shifted with each iteration,
we do not need to plan as meticulously as we would for a waterfall approach.

We then enter the section with the prototyping iterations, 
where we will alternate between the implementation of what we planned in the previous section and
the evaluation of the results of the previous iteration which serve as the basis for the next iteration.
 
Then a final evaluation of the whole project will be done, where the final results will be evaluated and the project will be concluded,
and advice for the pickup of the project will be given.

\section{Decision-Making}
\label{decision_making}

As previously described, the methodology of Prototyping benefits from a very tight loop of iterations between the different phases of the project.
While this is highly effective in producing a good end result, it can also take many iterations and a lot of experimentation until an adequate tool or solution has been found.
Given the constraints of a limited time frame for this project, it becomes crucial to use this time as efficiently as possible.
Sometimes, when the time does not permit a thorough exploration of 

To ensure that the decisions made are the most optimal within the constraints of the available information, adopting a systematic, replicable, and transparent decision-making process becomes essential. Over the years, various frameworks have been crafted to guide decision-making, 
particularly when information is complex and multidimensional.

\subsection{Weighted Sum Model}
Evangelos Triantaphyllou suggests that the \ac{WSM} is in practice the most used and most relevant decision-making framework \footcite[p. 1]{triantaphyllouIntroductionMultiCriteriaDecision2000}.
The \ac{WSM} method, by design, mandates the assignment of specific weights to each criterion based on its relevance. Subsequent to this, every alternative is evaluated based on these weighted criteria, resulting in a cumulative score.
The alternative with the highest score is therefore the optimal choice.


\begin{figure}[h]
    \centering
    \Large
    \[ A_{i}^{WSM-score} = \sum_{j=1}^{n} w_{j} a_{ij} \quad \ for \ i = 1, 2, 3, \dots, m. \]
    \caption{Formula for calculating the \ac*{WSM} score\protect\footnotemark} 
\end{figure}
\footnotetext{\cite{WeightedSumModel2022}}

Where:
\begin{itemize}
    \item \(w_{j}\): This represents the weight assigned to the \(j\)-th criterion. Weights are determined by the decision-makers based on the relative importance of each criterion. They should be normalized (i.e., the sum of all weights should be 1 or 100\%) to maintain a consistent scale.
    
    \item \(a_{ij}\): This represents the score or rating of the \(i\)-th alternative concerning the \(j\)-th criterion. This score is an assessment of how well the alternative meets or satisfies the specific criterion. 
    \end{itemize}

This method, despite its simplicity and direct approach, isn't without limitations.
One notable drawback is its dependence on dimensionless scales.
For the weights to properly reflect the criteria's importance, the scores need to be on a common, dimensionless scale, a detail not always feasible or convenient in practice.

\subsection{Simple Multi-Attribute Rating Technique}

In contrast to the \ac{WSM}, which predominantly utilizes a
direct mathematical approach to rank alternatives based on their weighted sum
scores, the \ac{SMART} methodology offers a more comprehensive approach to
multi-criteria decision-making. While \ac{WSM} is primarily concerned with simple
weighted arithmetic sums, the \ac{SMART} method dives deeper, ensuring that diverse
performance values—both quantitative and qualitative are harmonized and placed
on a common scale.

The \ac{SMART} method, grounded in \ac{MAUT}, provides a
structured framework that encompasses more than just the weighting of criteria.
It involves:

\begin{enumerate} \item Discernment of vital criteria pertinent to the decision
in focus. \item Weight allocation to each criterion in accordance to its
significance. \item Evaluation of each potential alternative against the
identified criteria, culminating in a score. \item Aggregation of these
individual scores via their associated weights, yielding a total score for every
alternative. \end{enumerate}

By adhering to the \ac{SMART} framework, alternatives can be sequenced based on
their aggregated weighted scores. This systematic approach equips
decision-makers to choose solutions that align closely with their objectives.
The computational formula integral to the \ac{SMART} method is:


\begin{figure}[h]
    \centering
    \Large
    \[ x_j = \frac{\sum_{i=1}^{m} w_i a_{ij}}{\sum_{i=1}^{m} w_i}, \quad j = 1, \ldots, n. \]
    \caption{Formula for calculating the \ac*{SMART} score\protect\footnotemark} 
\end{figure}

\footnotetext{Taken from \cite[p. 6]{fulopIntroductionDecisionMaking2005}}

Where:
\begin{itemize}
    \item \( x_j \) Is the overall utility score for alternative \( j \). The higher the score, the better the alternative, in comparison to the other alternatives.
    \item \( a_{ij} \)  Is the utility score for alternative \( j \) for the criterion \( i \).
    \item \( w_i \) Is the weight of criterion \( i \).
\end{itemize}

This method's emphasis on utility functions ensures a more nuanced and adaptable
approach to decision-making compared to models like WSM, making it suitable for
complex scenarios where criteria and alternatives are diverse in nature \footcite[p. 6]{fulopIntroductionDecisionMaking2005}.

\subsection{SMART Exploiting Ranks}
\label{smart_er}

The \ac{SMART-ER} method is a variant of the \ac{SMART} method that attempts to alleviate the largest issue \footcitetext[p. 26]{barfodMulticriteriaDecisionAnalysis2014} of the original \ac{SMART} method, 
namely the problem of a somewhat arbitrary ranking of the options if no numerical values can be derived.

This method addresses the issue by letting the decision maker simply ranking the different criteria in relation to each other and then normalizing the weights\footcite[p. 296]{robertsWeightApproximationsMultiattribute2002a}.
They propose the different weighting curves.

\begin{figure}[h]
    \centering
    \[ w_{i}(ROC) = \frac{1}{n} \sum_{j=1}^{n} \frac{1}{j}, \quad i = 1, \ldots, n. \]
    \caption{Formula for the \ac{ROC} weights}
\end{figure}

The \ac{ROC} takes the centroid of the rank order and uses the reciprocal of the rank as the weight.

\begin{figure}[h]
    \centering
    \[ w_{i}(\ RS\ ) = \frac{n + 1 - i}{n(n + 1)/2}, \quad i = 1, \ldots, n. \]
    \caption{Formula for the \ac{RS} weights}
\end{figure}

The \ac{RS} uses linear curve where weights are normalized by dividing them by the sum of all weights.

\begin{figure}[h]
    \centering
    \[ w_i(RR) = \left( \frac{\frac{1}{i}}{\sum_{j=1}^{n} \frac{1}{j}} \right) , \ rank\ i = 1, \ldots, n, \ option\ j = 1, \ldots, n \]
    \caption{Formula for the \ac{RR} weights}
\end{figure}

The \ac{RR} emphasizes the most important criteria by using the reciprocal of the rank as the weight, then normalizing each weight by the sum of all reciprocals.


\chapter{Conclusion}
\label{Conclusion}

\blindtext[2]
\blindtext[2]

% 5. Conclusion
% - Summary
% - Future Work


\chapter{Summary and Outlook}
\label{summary}
 
\section{Summary}

By way of this project paper we have assessed and described the state of the art for current technologies in the areas of 
containerized software, container orchestration and how those tie in with Software defined infrastructure.
We have done the same for the state of the art in the solution of complex problems found in the field of and  solved with \ac{HPC},
namely the two larger classes of problems, \ac{LCP} and \ac{TCP} problems, and what strategies are usually empoyed to solve those problems.

We have then addressed the problems identified by the problem statement with the statement of an initial goal, 
and structured an intervention to solve those problems by way of the creation of a prototype.
Before the iterative process of prototyping could begin we had to make a non-trivial decision on the choice of a container orchestrator/ workflow manager.
For which we defined, discussed and weighted the selection criteria, and then evaluated the available options against those criteria and 
We found our best fit through the employment of the \ac{SMART-ER} method, which was the \ac{k8s}-based orchestrator named Pachyderm.

The iterative process of prototyping was split into three main areas of focus, the infrastructure, the solution to \acp{TCP} and the integration of a complete 
\ac{CI/CD} pipeline. Each of these areas was given at least two iterations and the results of each iteration were documented, evaluated and recomendations for
future iterations were made.

The final results of the prototype were then evaluated against the initial goal and the problems identified in the problem statement,
and the resulting artifact was found to be a valid form of intervention to solve the problems identified in the problem statement,
while still being limited in scope and applicability, as was to be expected from a non production prototype. 


\section{Outlook}

% The prototype created in this project paper is a valid intervention to solve the problems identified in the problem statement,
% but still proves to have great potential for future work.
% As the time was limited and mostly spend on the furthering of the project, many avenues of high interest are still open for exploration.

% The two main areas of focus for future work are 
% the optimization of the communication stack on every level of abstraction,
% the expansion of usability functions relating to the Catalogue as these will bring the most immediate value to potential users





% ----------------------- Ende Inhalt ---------------------------------

%%% Aknowledgements %%%
% \clearpage
% \chapter*{Acknowledgements}
\label{Acknowledgements}

I would like to extend my sincerest gratitude to my supervisors, [Name] and [Name], 
for their invaluable support and guidance throughout the course of this research project.
Their insightful feedback and constructive criticism have been instrumental in shaping this paper.

I would also like to thank my friends and fellow students, Pascal Daniel und Max Kirschmann, 
for their contributions and insightful discussions that helped shape this research paper.

Additionally, I would like to acknowledge the contribution of GitHub Copilot, 
an AI-powered programming tool that provided assistance with the formulation of certain sections of this paper.
Its capabilities in generating relevant code snippets and suggestions were helpful in the completion of this project.

%%% Quellenverzeichnisse (keine Anpassung nötig) %%%
\clearpage
\literaturverzeichnis
%%% Ende Quellenverzeichnisse %%%


%%% Erklärung (keine Anpassungen nötig) %%%
% steht ganz am Ende des Dokuments
\cleardoublepage
\end{document}